\chapter{Resultados e Discussões}
\label{sec:resultados}

\section{Avaliação Preditiva e Eficiência Computacional}

A avaliação experimental do sistema modernizado foi conduzida utilizando os dados espectrométricos de cinco analitos (Mn, Mo, Cr, Ni e Fe) presentes em amostras de aço-ligas. Os modelos de Regressão Linear Múltipla (RLM), RLM-\textit{bagging} e RLM-\textit{subagging} foram comparados quanto à acurácia preditiva, mensurada pelas métricas RMSEP, MAE e $R^2$, e quanto ao tempo de processamento computacional.

Para o Molibdênio (Mo) e o Cromo (Cr), a RLM clássica obteve menor RMSEP, indicando que a introdução das estratégias de \textit{ensemble} não proporcionou ganho de acurácia preditiva no conjunto de teste avaliado. Os resultados detalhados demonstram que a estrutura linear simples da RLM foi suficiente para capturar a variabilidade desses elementos sem introduzir a variância adicional inerente ao processo de reamostragem.

No caso do Ferro (Fe), o modelo RLM-\textit{bagging} apresentou o menor erro de predição, enquanto para o Manganês (Mn) e o Níquel (Ni), o RLM-\textit{subagging} obteve desempenho superior ou equivalente às demais abordagens. O desempenho diferenciado entre os analitos reflete suas características espectrais individuais no conjunto de dados. A verificação detalhada das causas físico-químicas subjacentes e a decomposição formal de viés e variância constituem tópicos sugeridos para investigações futuras.

Em termos de eficiência computacional, o RLM-\textit{subagging} demonstrou tempos de execução sistematicamente inferiores aos do RLM-\textit{bagging} em quatro dos cinco analitos avaliados. Essa redução do tempo de processamento é atribuída ao menor tamanho da amostra utilizada em cada submodelo durante o procedimento de \textit{subagging} ($k < 1$), o que diminui o custo computacional do ajuste de matrizes por Mínimos Quadrados Ordinários (MQO).